%!TEX root = ../paper.tex
%
% The decidable path fragment.
%
% Companion to prototype/path_domain.cpp, which implements everything below and
% validates it on the paths appearing in Figure \ref{fig:rculist}, Section
% \ref{sec:eval}, and Appendix \ref{appendix:bst_del}.
%
\section{A Decidable Path Fragment}
\label{sec:pathfrag}

\newcommand{\pvar}{\pi}
\newcommand{\MayAlias}{\textsf{MayAlias}}
\newcommand{\defn}{\;\triangleq\;}
\newcommand{\Fields}{\mathcal{F}}

The type system's framing side conditions are all negated \MayAlias{} queries,
and \textsc{T-ReIndex} is applied by hand throughout the paper.  Both must become
algorithms before the system can be checked mechanically.  This section fixes a
fragment of the path language on which they are, states the algorithms, and
records what falls outside.

\subsection{A gap in the notation: length indices versus word sharing}
\label{sec:pathfrag-gap}

The paper writes an abstracted traversal as $(f_1|\ldots|f_n)^{k}$ and describes
$k$ as ``the number of loop iterations run''.  Read literally, $\rho$ constrains
only the \emph{length} of the traversal: $(\textit{Left}|\textit{Right})^{k}$
denotes the set of all words of length $k$ over $\{\textit{Left},
\textit{Right}\}$, and two types mentioning $(\textit{Left}|\textit{Right})^{k}$
are related only by agreeing on $k$.

That reading is too weak for the invariants the paper relies on.  Consider the
BST loop, at the point where the outer traversal has found its node:
\[
\textit{pr} : \textsf{rcuItr}\;(l|r)^{k}\;\{l|r \rightharpoonup \textit{cr}\}
\qquad
\textit{cr} : \textsf{rcuItr}\;(l|r)^{k}.(l|r)\;\{\}
\]
The field map on \textit{pr} asserts that \textit{cr} is \emph{its} child.  Under
the length reading, \textit{pr} could sit at $l.l$ and \textit{cr} at $r.r.l$ ---
both consistent with $k = 2$ --- and the field map would not hold.  What the
proofs actually need is that the two occurrences of $(l|r)^{k}$ denote the
\emph{same word}, so that \textit{cr}'s path extends \textit{pr}'s.

The same gap shows up in Section \ref{sec:eval}, which describes $T_0$'s
traversal as ``summarised as $(l|r)^{k}$'' --- summarising \emph{one} traversal,
i.e.\ one word.  We therefore make the sharing explicit rather than leaving it to
the reader.

\paragraph{Consequences.}  This is not only a presentation matter.  Under the
length reading, the field-map conjunct of $\llbracket \textsf{rcuItr}\,\rho\,
\mathcal{N} \rrbracket$ is not implied by the path conjunct, and the
\textsc{Unlink} and \textsc{Replace} lemmas' appeals to ``proper linkage'' do not
go through.  Any mechanization has to pick a reading; it should pick this one.

\subsection{Syntax}
\label{sec:pathfrag-syntax}

Let $\Fields$ be the finite set of \RCU-typed field names, which in a C program
comes from the struct declaration and is small.

\[
\begin{array}{r@{\;}c@{\;}l@{\qquad}l}
\rho &::=& \epsilon \mid \rho \cdot s & \text{paths} \\
s    &::=& f                 & \text{one step along } f \in \Fields \\
     &\mid& D                & \text{one step along some } f \in D \subseteq \Fields \\
     &\mid& \pvar : D        & \text{the word bound to } \pvar,\ \text{which lies in } D^{*}
\end{array}
\]

A path variable $\pvar$ replaces the paper's $(f_1|\ldots|f_n)^{k}$.  Variables
are scoped at the type environment, not the type: every occurrence of $\pvar$
anywhere in $\Gamma$ denotes the same word.  For the same reason its alphabet
$D$ is a property of the \emph{variable} and not of an occurrence, so a
well-formed environment assigns each $\pvar$ one $D$ throughout; occurrences
carrying different alphabets denote different variables and must be renamed
apart.  (This is not pedantry: violating it was the one thing that made our
property tests fail, and it silently breaks reindexing, which compares segments
up to their alphabets.)  A \emph{valuation} $\theta$ maps
each variable to a concrete word in its alphabet, and
\[
\llbracket \rho \rrbracket\theta \;\subseteq\; \Fields^{*}
\qquad
\llbracket f \rrbracket\theta = \{f\}
\quad
\llbracket D \rrbracket\theta = D
\quad
\llbracket \pvar \rrbracket\theta = \{\theta(\pvar)\}
\]
extended to concatenation pointwise.  Note $\llbracket \rho \rrbracket\theta$ is
finite: it is a product of at most $|\rho|$ finite sets.  A type environment
denotes $\exists\theta \ldotp \bigwedge_i \exists p_i \in \llbracket \rho_i
\rrbracket\theta \ldotp h^{*}(rt,p_i) = s(x_i,\mathit{tid})$ --- one valuation
shared across the environment, but each variable free to resolve its own
disjunctive steps.

\subsection{The fragment}
\label{sec:pathfrag-def}

\begin{definition}[Cursor Normal Form]
\label{def:cnf}
A path is in \emph{CNF} when no path variable occurs in it more than once.  A
type environment is in CNF when every path is, and the variables are drawn from a
fixed list $\pvar_1 < \cdots < \pvar_v$ mentioned in increasing order by every
path.
\end{definition}

$v$ is the number of simultaneously live cursors, which is the loop-nesting depth
of the traversal: $v = 1$ for the bag, $v = 2$ for the BST delete, whose two
loops are exactly the ``multiple cursors'' Section \ref{sec:eval} describes.  It
is bounded statically by the program text.

The restriction earns its keep by ruling out the two shapes that would force
general word equations: a variable occurring twice ($\pvar \cdot f \cdot \pvar$),
and misaligned variables ($\pvar \cdot a$ against $b \cdot \pvar'$, which would
require solving $\pvar \cdot a = b \cdot \pvar'$).  We reject these rather than
approximate them, so that a programmer can see from the syntax whether their code
is in scope.

\subsection{MayAlias}
\label{sec:pathfrag-alias}

\[
\MayAlias(\rho,\rho') \defn
  \exists\theta \ldotp \llbracket \rho \rrbracket\theta \cap \llbracket \rho' \rrbracket\theta \neq \emptyset
\]
By the corrected \textbf{UNQR} (\S\ref{sec:inv-revised}), distinct concrete paths
from the root reach distinct nodes, so two references may alias exactly when
their path languages can meet.

The decision procedure has two stages.

\paragraph{Stage 1 --- consume the maximal identical prefix.}
While $\rho$ and $\rho'$ agree segment for segment, drop the common segment.
This is \emph{exact}: aliasing requires the two whole words to be equal, so the
prefixes must resolve identically, and a shared variable in the prefix therefore
contributes the same word to both sides.

\paragraph{Stage 2 --- relax and intersect.}
On the remainders, replace each variable $\pvar : D$ by its regular constraint
$D^{*}$, build the obvious NFA for each path, and test the product for a
reachable accepting state.

\begin{lemma}[Soundness]
\label{lem:alias-sound}
If the procedure answers \textsf{no}, then $\neg\MayAlias(\rho,\rho')$.
\end{lemma}
\begin{proof}[Proof sketch]
Stage 1 preserves the set of solutions exactly.  Stage 2 discards the constraint
that two occurrences of the same variable agree, which can only enlarge the
languages being intersected; so a \emph{no} on the relaxed problem is a \emph{no}
on the original.
\end{proof}

Soundness in this direction is what the type system needs: every use of
\MayAlias{} in \textsc{T-UnlinkH}, \textsc{T-Replace} and \textsc{T-Insert} is
under a negation, so over-approximating aliasing makes the checker reject safe
programs but never accept unsafe ones.

The procedure is not complete --- after the prefix diverges, sharing is
forgotten.  It is exact whenever the shared variables align, which covers every
path in the paper; \S\ref{sec:pathfrag-eval} lists what was checked.

Stage 1 is linear.  Stage 2 is a product construction over
$(|\rho|+1)\times(|\rho'|+1)$ states with $|\Fields|$ letters, so the whole
procedure is $O(|\rho|\cdot|\rho'|\cdot|\Fields|)$.

\subsection{Join}
\label{sec:pathfrag-join}

Merging two branches computes a segment-wise least upper bound, which is exactly
\textsc{T-PSub1}/\textsc{T-PSub2} ($\rho.f_1 \prec: \rho.f_1|f_2$): concrete and
disjunctive steps join by union of their field sets; a variable joins only with
itself.  When the shapes do not line up the join is undefined and the type system
falls back to \textsc{T-TSub} ($\textsf{rcuItr}\,\_ \prec: \textsf{undef}$),
which is sound and merely imprecise.  This is the operation that turns Figure
\ref{fig:altfield}'s two branches into $\textit{Left}|\textit{Right}$.

\subsection{Reindexing}
\label{sec:pathfrag-reindex}

\textsc{T-ReIndex} closes a loop's back edge.  Automating it is first-order
matching: find a single substitution $\sigma$ with $\sigma(\Gamma_{\mathit{entry}})
= \Gamma_{\mathit{exit}}$.  Restricting $\sigma$ to suffix extension,
$\sigma(\pvar) = \pvar \cdot w$, covers every loop in the paper and makes
matching linear --- locate the variable, read off the spliced segments, then
verify the candidate against the rest of the environment.

The verification step is where the paper's side condition lands.  Section
\ref{subsection:typwrt} argues reindexing is sound because ``either all are
reindexed, or none''; requiring one $\sigma$ to validate the whole environment
is precisely that condition, discharged by construction.

For the bag this yields $\sigma(k) = k.\textit{Next}$, mapping lines 19--20 of
Figure \ref{fig:rculist} to lines 23--24.  For the BST inner loop it yields
$\sigma(m) = m.\textit{Left}$ and leaves $k$ untouched --- the two cursors are
reindexed independently, which is what lets the two unlinks at different depths
be framed apart.

\subsection{Widening}
\label{sec:pathfrag-widen}

Reindexing closes a back edge only once a variable exists to extend.  On first
arrival there is none --- the bag enters its loop with
$\textit{par} : \epsilon$ and $\textit{cur} : \textit{Next}$ --- so the variable
must be \emph{introduced}.  This is also what makes path length bounded: it
grows without limit along a loop, so the domain has no finite height and
reindexing alone gives no termination argument.

The essential constraint is that \textbf{widening cannot be done path by path.}
Generalising each path at its own longest common prefix turns the bag into
\[
\textit{par} : (\textit{Next})^{k}
\qquad
\textit{cur} : \textit{Next}.(\textit{Next})^{k}
\]
which denotes the right \emph{language} but asserts that \textit{cur} is
\textit{par} with a step \emph{prepended}.  The relational structure ---
\textit{cur} is \textit{par}'s child --- is destroyed, and with it every field
map that depends on it.  Over a single-field alphabet the two forms happen to
coincide, so the bag alone would not expose the error; the BST, where
$(l|r).(l|r)^{k}$ and $(l|r)^{k}.(l|r)$ differ, does.

So a single fresh variable is inserted at one common position across the whole
environment.  Writing $\Gamma(x)$ for $x$'s path, find the smallest $i$ and a
single chunk $w$ with
\[
\Gamma_{\mathit{exit}}(x) \;=\; \Gamma_{\mathit{entry}}(x)[0,i) \cdot w \cdot \Gamma_{\mathit{entry}}(x)[i,\cdot)
\qquad\text{for every \emph{moving} } x
\]
and replace $w$ by a fresh $\pvar : D$, where $D$ is the union of the alphabets
of $w$'s segments.  Paths the loop leaves alone --- those with
$\Gamma_{\mathit{entry}}(x) = \Gamma_{\mathit{exit}}(x)$, such as the outer
cursors during the BST's inner loop --- pass through untouched.  Cursors that
advance by different amounts are refused, as outside the fragment.

\begin{lemma}[Widening is an upper bound]
\label{lem:widen-sound}
$\theta(\pvar) = \epsilon$ recovers $\Gamma_{\mathit{entry}}$ exactly, and
$\theta(\pvar)$ ranging over the words of $w$ covers $\Gamma_{\mathit{exit}}$.
\end{lemma}
\begin{proof}[Proof sketch]
Immediate from the displayed equation, since all moving paths share one $w$, so
a single valuation serves the whole environment simultaneously.
\end{proof}

This is how one \emph{infers} $(\textit{Next})^{k}$ rather than being handed it.
The loop analysis is then: apply the body; if the edge closes by equality or by
reindexing, that is the invariant; otherwise widen and retry, to a fixed bound.

\subsection{What this does and does not settle}
\label{sec:pathfrag-decid}

It does not make the declarative type system decidable; the concession in the
authors' response stands, because a declarative derivation may insert
\textsc{T-Conseq} anywhere.  What it gives is a decidable \emph{algorithmic}
presentation: with a canonical join at merges, reindexing-by-matching at back
edges, and \MayAlias{} decided as above, subsumption is applied only at fixed
points, and checking a critical section becomes a terminating forward dataflow
over CNF environments.  The fragment is the class of programs that algorithm is
complete for.

This distinction is the honest answer to Reviewer 2's decidability objection and
to Reviewer 3's automation question, and it is the property the Clang frontend
needs: \textsf{Sema} has no symbolic execution to fall back on, so the check must
be deterministic and terminating on its own.

\subsection{Status}
\label{sec:pathfrag-eval}

\texttt{prototype/path\_domain.h} implements CNF membership, \MayAlias{}, join,
substitution, reindexing, widening and the loop fixpoint in dependency-free
C++17, structured to be lifted into \textsf{Sema}.
\texttt{prototype/path\_domain\_test.cpp} passes 30/30.

\paragraph{Unit tests} use only paths that appear in the paper.
\begin{itemize}
\item \textbf{Bag.}  $\textit{par} : (\textit{Next})^{k}$ and $\textit{cur} :
  (\textit{Next})^{k}.\textit{Next}$ are correctly found non-aliasing --- the
  case a plain regular-language intersection gets wrong, and the one Stage 1
  exists for.
\item \textbf{BST.}  The two-cursor environment is in CNF; the cursors are found
  non-aliasing, both with each other and with the outer cursor \textit{cr}.
\item \textbf{Join.}  $\textit{Left}.\textit{Left} \sqcup
  \textit{Right}.\textit{Right} = (\textit{Left}|\textit{Right}).(\textit{Left}|\textit{Right})$,
  reproducing Figure \ref{fig:altfield}'s merge.
\item \textbf{Invariant inference.}  Starting from the environment on entry to
  each loop and applying only the body, the analysis \emph{derives} the paper's
  hand-written invariants.  The bag converges in one widening to exactly
  $(\textit{Next})^{k}$ and $(\textit{Next})^{k}.\textit{Next}$ --- lines 19--20
  of Figure \ref{fig:rculist}.  The BST outer loop converges to $(l|r)^{k}$ and
  $(l|r)^{k}.(l|r)$.  The BST inner loop converges with the outer cursors left
  frozen, reindexing $m$ alone, and the inferred invariant still separates the
  two cursors --- the property the two unlinks at different depths depend on.
\item \textbf{Boundary.}  $\pvar . l . \pvar$ is rejected as outside CNF, and
  unevenly advancing cursors are refused rather than approximated.
\end{itemize}

\paragraph{Property tests} check the abstract operations against the concrete
semantics by enumerating valuations up to a length bound.  Over $2\times10^{5}$
random path pairs across five seeds:
\begin{itemize}
\item \textbf{P1} \MayAlias{} never misses a real alias --- the soundness
  direction of Lemma \ref{lem:alias-sound}.  Roughly $31\%$ of random pairs do
  alias, so the test is not vacuous.  Measured imprecision is $4.2\%$ of pairs
  answered conservatively \textsf{yes}; raising the enumeration bound from $3$
  to $5$ reclassifies about $6\%$ of those as genuine aliases, so the residue is
  Stage 2 forgetting sharing.
\item \textbf{P2} join over-approximates both operands.
\item \textbf{P3} widening is an upper bound, in both directions of Lemma
  \ref{lem:widen-sound}.
\item \textbf{P4} reindexing recovers any suffix-extension substitution.
\end{itemize}

\paragraph{Outstanding.}  A completeness characterisation for Stage 2: we claim
only soundness and exactness-when-aligned, and the measured $4.2\%$ is an upper
bound on what a sharing-aware Stage 2 could recover.  Whether CNF alignment
makes Stage 2 exact outright is open, and worth settling before the imprecision
shows up as a false positive in a compiler diagnostic.
